\chapter{Robustness of LQR Solvers in Different Dimensions}
\label{chap:lqr}

\section{A Known Model}
\label{sec:known-model}

We start by comparing two approaches that have knowledge of the underlying model,
and thus do not need to estimate gradients but can calculate them exactly. This allows
us to clearly separate the benefits of a low-dimensional controller from estimation
errors.

\subsection{Setting}
\label{subsec:known-model-setting}

% mb pg
Fazel et al.~\cite{Fazel2018Global} introduced global convergence guarantees for policy gradient methods based on the
true model in 2019. The paper describes a method for calculating the gradient of a policy from
the system and cost matrices, and updating the policy accordingly. As the paper uses the same $u_k = -K x_k$
convention as we do, we can implement the loop directly in the following way.

% Def mb_pg
Starting from an initial, stabilizing gain $K_0 \in \mathcal{K}$ and the second moment $\Sigma_0$ of the initial
state distribution, every iteration $j$ first solves the discrete Lyapunov equation \eqref{eq:lyapunov}
for $P_{K_j}$ using the implementation from SciPy~\needcite{scipy}. The state correlation matrix
$\Sigma_{K_j}$ then follows from solving a second Lyapunov equation,
\begin{equation}
    \label{eq:state-correlation}
    \Sigma_{K_j} = (A - B K_j) \Sigma_{K_j} (A - B K_j)\T + \Sigma_0 ,
\end{equation}
where, for initial states drawn uniformly from $[-10, 10] \times [-1, 1]$ as described in \cref{subsubsec:lqr-cost-of-gain},
$\Sigma_0 = \operatorname{diag}\left( 10^2/3, \, 1^2/3 \right)$. The policy gradient
\begin{equation}
    \label{eq:pg-gradient}
    \nabla f(K_j) = 2 \left( \left( R + B\T P_{K_j} B \right) K_j - B\T P_{K_j} A \right) \Sigma_{K_j}
\end{equation}
is then used for a gradient step on the gain matrix. \Cref{alg:mb-pg} summarizes the loop.
\begin{algorithm}[htbp]
    \caption{Model-based PG}
    \label{alg:mb-pg}
    \begin{algorithmic}[1]
        \REQUIRE stabilizing gain $K_0 \in \mathcal{K}$, second moment $\Sigma_0$, step size $\eta$
        \FOR{$j = 0, 1, 2, \dots$}
            \STATE solve \eqref{eq:lyapunov} for $P_{K_j}$
            \STATE solve \eqref{eq:state-correlation} for $\Sigma_{K_j}$
            \STATE compute the policy gradient $\nabla f(K_j)$ by \eqref{eq:pg-gradient}
            \STATE update the gain by $K_{j+1} = K_j - \eta \nabla f(K_j)$
        \ENDFOR
    \end{algorithmic}
\end{algorithm}

The fixed learning rate $\eta$ is a hyperparameter, which we set to $\eta = 10^{-5}$ for all inflation sizes, based on
experiments whose results can be found in \cref{sec:app-learning-rate-mb-pg}. This
policy gradient approach serves us as the basic, fully-parameterized implementation for this experiment.

% mb spi

Now we also want to implement a solver that produces a low-dimensional controller. As in our
case the controller is characterized by the gain matrix, and for matrices dimensionality is ``often
[\ldots] expressed by means of the rank''~\cite{Recht2010Guaranteed}, we aim at learning a low-rank
gain matrix. However, to make this comparable later to RL approaches that rely on the number of parameters,
we measure the complexity of a gain by its policy size $\lvert K \rvert$ from \cref{def:policy-size}, which uses
the more efficient version between the full and the factored parameterization.

Recht et al.~\cite{Recht2010Guaranteed} have already shown that nuclear norm minimization introduces
low-rank solutions for convex matrix problems. However, as we are dealing with a non-convex optimization
problem~\cite{Fazel2018Global}, we cannot transfer the results over directly. Here, Park et al.~\cite{Park2020Structured}
offer us an approach to impose different structures on our policy during learning via regularizers
based on norms, one of them being the nuclear norm. We implement the model-based approach introduced there
with small adaptations in the following way.

% Def mb_spi
We reuse the computation of the gradient $\nabla f(K_j)$ from \cref{alg:mb-pg}.
For inducing the structure, we use an adaptive learning rate combined with a line search criterion.
Starting from the initial step size $\eta = 1 / (2 \lambda)$, each iteration takes a gradient step and applies the
proximal operator of the nuclear norm, which soft-thresholds the singular values,
\begin{equation}
    \label{eq:spi-prox}
    K^{+} = U \operatorname{diag} \left( \max\{\sigma_1 - \eta \lambda, 0\}, \dots,
    \max\{\sigma_{\bar{r}} - \eta \lambda, 0\} \right) V\T ,
    \qquad \bar{r} \coloneqq \min\{d_{\mathrm{out}}, d_{\mathrm{in}}\} ,
\end{equation}
where $U \operatorname{diag}(\sigma_1, \dots, \sigma_{\bar{r}}) V\T$ is the singular value decomposition of
the next policy candidate $K^{+} = K_j - \eta \nabla f(K_j)$. The candidate $K^{+}$ is accepted if it is stabilizing, $\spr(A - B K^{+}) < 1$, and
satisfies the sufficient decrease condition
\begin{equation}
    \label{eq:spi-sufficient-decrease}
    f(K^{+}) \leq f(K_j) - \eta \tr \left( \nabla f(K_j)\T G \right) + \frac{\eta}{2} \lVert G \rVert_F^2 ,
    \qquad G \coloneqq \frac{1}{\eta} \left( K_j - K^{+} \right) .
\end{equation}
Otherwise, the step size is reduced to $\beta \eta$ and the proximal step is repeated. The learning loop further can be
stopped earlier if the criterion $\lVert K_{j+1} - K_j \rVert_F \leq \varepsilon$ is satisfied.
\Cref{alg:mb-spi} summarizes the loop.
\begin{algorithm}[htbp]
    \caption{S-PI}
    \label{alg:mb-spi}
    \begin{algorithmic}[1]
        \REQUIRE stabilizing gain $K_0 \in \mathcal{K}$, second moment $\Sigma_0$, weight $\lambda$, factor $\beta$,
        tolerance $\varepsilon$
        \FOR{$j = 0, 1, 2, \dots$}
            \STATE compute $\nabla f(K_j)$ as in \cref{alg:mb-pg}
            \STATE set the initial step size $\eta = 1 / (2 \lambda)$
            \FOR{at most $20$ line search steps}
                \STATE compute the candidate $K^{+}$ by \eqref{eq:spi-prox}
                \IF{$\spr(A - B K^{+}) < 1$ and \eqref{eq:spi-sufficient-decrease} holds}
                    \STATE \textbf{break}
                \ENDIF
                \STATE reduce the step size, $\eta \leftarrow \beta \eta$
            \ENDFOR
            \STATE update the gain, $K_{j+1} = K^{+}$
            \IF{$\lVert K_{j+1} - K_j \rVert_F \leq \varepsilon$}
                \STATE \textbf{stop}
            \ENDIF
        \ENDFOR
    \end{algorithmic}
\end{algorithm}

As one can see, there are a lot of parameters that can be modified here. We set $\varepsilon = 10^{-8}$,
and use the same $\beta = 0.5$ as in the experiments section of Park et al.~\cite{Park2020Structured}.
For $\lambda$, as it defines the learning rate as well, we have run an experiment with various values,
which can be found in \cref{sec:app-spi-lambda}. A value of $\lambda = 5$ has shown the most promising results across the sizes.

One thing that we implemented differently than the paper is the cap on the number of line search iterations in line four, which we set to 20.
Some of our previous experiments had shown that if we leave it uncapped, for unstable updates which are impossible
to stabilize, the loop runs until the update produces no effect at all.

Now that we have defined our controllers and how they learn, we test them in comparison to the DARE-calculated
solution \eqref{eq:dare} from \cref{subsubsec:lqr-optimal}.

\subsection{Results}
\label{subsec:known-model-results}

We apply our solvers to the true models base and inflated versions, for now noise-free, as described in \needcite{cref inflation loop}. 
Specifically, as in \needcite{cref to test protocol}, we train 10 controllers each per solver, before evaluating them on the test episodes.
The results are then pooled.

The most interesting finding is, how much faster SPI converges compared to the Policy Gradient. While it achieves its best training
suboptimality within around 50 iterations, MB PG still falls drastically short of it after 450 iterations. We can also see in 
\needcite{table showing rank / complexity} that the reduction in complexity for all inflation sizes has worked well with SPI. 
For our test suboptimality, SPI also outperforms MB PG here by at least one order of magnitude. \needcite{across sizes suboptimality}
visualizes this.

% the table and figure from above

Overall, this looks very promising for SPI, still, our goal is to find out whether it also is a robust controller.

\subsection{Noise Experiments}
\label{sec:true-model-noise}

Now, we reuse the controllers we trained in \needcite{cref to subsec above} and evaluate them under the test protocol
with noise\needcite{cref to noise introduction} injected. 

When looking at the results, one can notice that our baseline controller, which was optimal for the clean loop, mostly
performs worse than our trained controllers, except for external disturbance. However, when only comparing MB PG with
SPI, we can see differences in their performance, depending on the type of noise that is applied, as well as whether the
loop is inflated or not. For parametric noise for example, MB PG performs on average better than SPI here, however also
has a higher variance. When looking at external disturbance under large inflation, one can notice that with increasing
noise it becomes clearer that SPI is more robust here compared to MB PG. While for weak and standard their total spread
still shows some overlap, for strong and divergent SPI shows a cleanly seperated suboptimality below MB PG.

% figure of comparison across the sizes for external disturbance

For control noise in comparison, we can see that here MB PG is more robust on average across settings. However, its worst case
performance here is worse than SPI, so it has an overall wider spread of suboptimalities. Still, both mostly outperform the
baseline controller. Measurement noise becomes an interesting case here. While both controllers outperform the baseline, under
no inflation, MB PG seems to be much more robust here. However, when looking at the controllers under large inflation with the
same noise level applied, suddenly SPI shows more robustness here with little variation. MB PG on the other hand is only on
average better than the baseline, but in some cases worse. \needcite{figure below} visualizes this.

% left base noise measurement divergent, right large noise measurement divergent

Overall, for the strong level and large inflation, this is how both controllers compare to the baseline and each other across the
noise channels \needcite{table below}.

% table true model noise large channels

\subsection{Discussion}
\label{sec:true-model-discussion}

Overall, we note that trained controllers are mostly at an advantage compared to statically calculated controllers in regards to
robustness. When comparing now trained fully parameterized and structured approaches, differences definitely appear. Changes in
the underlying model, as in Parametric Uncertainty, affect all controllers mostly the same, even though a fully parameterized
version might offer better robustness on average. Looking at the other cases, we find that low-dimensional controllers actually
offer better robustness in noise related to the inputs of the controller, especially the larger the input dimension. Measurement
noise shows this clearly, as the performance of a overparameterized controller degrades here with increasing size, while a
regularized controller stays mostly the same, or even shows improvement. However, related to the output dimensions, under control
noise reduced controllers perform worse on average, however still better in the worst case than the full controller. And under large
size the reduction in needed parameters is of two magnitudes of order, which shows the efficiency of the reduced controller.

\section{Unknown Model across Dimensions}
\label{sec:unknown-model}

Now we take a look at the more realistic case, where we do not know the underlying model, and compare different approaches.
Before we can get into the results, we first take a deeper look at the solvers that we use and how the experiments are
set up.

\subsection{Setting}
\label{subsec:unknown-model-setting}

First, we introduce our basis for the indirect model-free approach, which is based on system identification. N4SID,
introduced by Van Overschee and De Moor~\cite{VanOverschee1994N4SID}, uses ``the projection of the future outputs [of the system]
onto the past and future inputs and the past outputs''~\cite[Sec.~3]{VanOverschee1994N4SID}. These projections are then called
``a function of the system matrices and the input-output block Hankel matrices''~\cite[Sec.~3]{VanOverschee1994N4SID}, which
allows us to reconstruct exactly that from the results. Important are further the assumptions that the systems are
observable (\cref{def:observability}) and controllable (\cref{def:controllability}). For simplicity, we rely on the
implementation of N4SID from SIPPY~\needcite{TODO}, which we integrate into our own setup. This implementation allows us
to recover the system matrices $A$, $B$ and $C$, whereas we implemented the recovery of the cost matrices $Q$ and $R$
from scratch based on~\cite{VanOverschee1994N4SID}.

The collection of observation and action data runs the plant for $T$ steps under the initial gain, with a randomly
sampled excitation added on top of it,
\begin{equation}
    \label{eq:identification-input}
    v_k = -K_0 y_k + \sigma_{\mathrm{e}} w_k , \qquad w_k \sim \mathcal{N}(0, I_{d_{\mathrm{out}}}) ,
    \qquad k = 0, \dots, T - 1 .
\end{equation}

As mentioned in the paper, ``as $i$ grows larger, the approximation error grows smaller''~\cite[Sec.~5.3]{VanOverschee1994N4SID},
with $i$ being dependent on the length of our trajectory described above used for identification. The other variable $j$
defines the dimensions of the Hankel matrices, and thus how far the projections reach into the future and the past.
As we want to be computationally efficient while still achieving the best performance, we evaluated different values
for both, to find the optimal setting to use in our experiment. The full results of the selection can be found in the
\cref{sec:app-identification}. Specifically, we use the settings in
\cref{tab:identification-settings}.
\begin{table}[htbp]
    \centering
    \caption{Identification settings per inflation size.}
    \label{tab:identification-settings}
    \begin{tabular}{@{}lrr@{}}
        \toprule
        Name & Trajectory length $T$ & Hankel horizon \\
        \midrule
        \emph{base}   & 500  & 5 \\
        \emph{small}  & 500  & 5 \\
        \emph{medium} & 500  & 5 \\
        \emph{large}  & 8000 & 5 \\
        \bottomrule
    \end{tabular}
\end{table}

The resulting identified system then serves as a surrogate for the true model
for the model-based approaches from \cref{subsec:known-model-setting}. This thus produces results without
ever having access to the real model.

An attentive reader might now ask why we would even want to use these iterative approaches, if we could just
solve the DARE \eqref{eq:dare} directly on this estimation. And indeed we include it in our following comparison.
However, to make this work with our construction of inflated state models, we need some modification to N4SID.
This stems from the fact that we need to calculate $K = \left( R + B\T P B \right)^{-1} B\T P A$ from
\eqref{eq:optimal-gain}, hence $R + B\T P B$ needs to be invertible.
But for our inflated space, the estimated $B$ and $R$ have the same null space, making the matrix singular on that subspace.
So to be able to apply this solver, we eliminate the additional singular values of $B$ that exceed the estimated order
of the model, thus reducing the model in dimension. This way we do not access the model, but create encoding and
decoding matrices that map between our dimensionally reduced model and the actual state observation, as well as the
actuation output. For the exact implementation, please refer to the repository~\needcite{our repository}.

Additionally, we also use direct model-free approaches, which are variations of the model-based approaches
from \cref{sec:known-model}. Unlike before, they do not calculate the gradient directly, but estimate
it based on the returns of $N_{\mathrm{traj}}$ slightly perturbed versions of the current policy, simulated for $l$ steps. As one might
assume, more ``data'', could conversely to better estimation, and thus better policy updates. So, to
settle this, we evaluated different combinations of them in \cref{sec:app-sample-size}. The best
trade-off between computational efficiency and quality gave $N_{\mathrm{traj}} = 100$ and $l = 800$.

% mf npg
We use the model-free natural policy gradient (NPG) estimation described by Fazel et al.~\cite{Fazel2018Global},
which at its core consists of the loop in \cref{alg:mf-gradient}. Instead of the proposed single-point estimator,
we use a two-point estimator inspired by Duchi et al. and Mania et al.~\cite{Duchi2015Optimal,Mania2018Simple},
as this drastically reduced the variance of the estimates, and thus the instability of the controller.

\begin{algorithm}[htbp]
    \caption{Two-point gradient estimation}
    \label{alg:mf-gradient}
    \begin{algorithmic}[1]
        \REQUIRE gain $K \in \R^{d_{\mathrm{out}} \times d_{\mathrm{in}}}$, number of perturbation pairs $N_{\mathrm{traj}}$,
        rollout length $l$, smoothing radius $r$
        \FOR{$i = 1, \dots, N_{\mathrm{traj}}$}
            \STATE sample $\Xi_i$ uniformly from the matrices with $\lVert \Xi_i \rVert_F = r$
            \STATE sample an initial state $x_0 \sim \mathcal{D}$
            \STATE simulate $K + \Xi_i$ and $K - \Xi_i$ for $l$ steps from $x_0$, with stage costs $c_k^{\pm}$ and
            observations $y_k^{\pm}$
            \STATE $\hat{J}_i^{\pm} = \sum_{k=0}^{l-1} c_k^{\pm}$, \quad
            $\hat{\Sigma}_i^{\pm} = \sum_{k=0}^{l-1} y_k^{\pm} \left( y_k^{\pm} \right)\T$, \quad
            $\delta_i = \hat{J}_i^{+} - \hat{J}_i^{-}$
        \ENDFOR
        \STATE $\bar{\delta} = \frac{1}{N_{\mathrm{traj}}} \sum_{i=1}^{N_{\mathrm{traj}}} \delta_i$
        \STATE $\widehat{\nabla f}(K) = \frac{1}{N_{\mathrm{traj}}} \sum_{i=1}^{N_{\mathrm{traj}}}
        \frac{d_{\mathrm{out}} d_{\mathrm{in}}}{2 r^2} \left( \delta_i - \bar{\delta} \right) \Xi_i$
        \STATE $\hat{\Sigma}_K = \frac{1}{N_{\mathrm{traj}}} \sum_{i=1}^{N_{\mathrm{traj}}} \left( \hat{\Sigma}_i^{+} + \hat{\Sigma}_i^{-} \right)$
        \STATE $\hat{f}(K) = \frac{1}{N_{\mathrm{traj}}} \sum_{i=1}^{N_{\mathrm{traj}}} \frac{1}{2} \left( \hat{J}_i^{+} + \hat{J}_i^{-} \right)$
        \RETURN $\widehat{\nabla f}(K)$, $\hat{\Sigma}_K$, $\hat{f}(K)$
    \end{algorithmic}
\end{algorithm}

We further aim to clarify a couple of things about the given loop.
Subtracting the mean difference $\bar{\delta}$ does not change the expected value of the estimate, since
$\E[\Xi_i] = 0$, but lowers its variance. A rollout whose observation exceeds $\lVert y_k \rVert = 10^{5}$,
we count as diverged. It is therefore ended early and charged a penalty cost of $10^{10}$. 
Additionally, for inflated dimensions, $\hat{\Sigma}_K$ is singular, so we precondition with its
truncated pseudoinverse $\hat{\Sigma}_K^{\dagger}$ instead, which discards singular values below $10^{-6}$ times the largest one.

Our full learning algorithm is then given in \cref{alg:mf-npg}.
\begin{algorithm}[htbp]
    \caption{Model-free NPG}
    \label{alg:mf-npg}
    \begin{algorithmic}[1]
        \REQUIRE stabilizing gain $K_0 \in \mathcal{K}$, $N_{\mathrm{traj}}$, $l$, $r$, step size $\eta$
        \FOR{$j = 0, 1, 2, \dots$}
            \STATE compute $\widehat{\nabla f}(K_j)$ and $\hat{\Sigma}_{K_j}$ by \cref{alg:mf-gradient}
            \STATE update the gain by $K_{j+1} = K_j - \eta \, \widehat{\nabla f}(K_j) \, \hat{\Sigma}_{K_j}^{\dagger}$
        \ENDFOR
    \end{algorithmic}
\end{algorithm}

After some evaluations, we have decided on a smoothing parameter of $r = 0.5$, and a learning rate of $\eta = 0.01$, based
on the results in \cref{sec:app-learning-rate-mf-npg}.

% spi mf
Next, we implement a model-free version of S-PI that relies on the same gradient
estimation (\cref{alg:mf-gradient}) as Model-free NPG. Another crucial difference,
apart from the estimation of the gradient compared to the calculation, is the ``usage of a fixed stepsize
$\eta$, rather than an adaptive stepsize from a backtracking linesearch''~\cite[Sec.~3.1]{Park2020Structured}, as that
would require access to the function value $f(K) = \tr(P_K \Sigma_0)$. Before we look at how we select this step size
and $\lambda$, we define our learning loop in \cref{alg:mf-spi}.
\begin{algorithm}[htbp]
    \caption{S-PI (model-free)}
    \label{alg:mf-spi}
    \begin{algorithmic}[1]
        \REQUIRE stabilizing gain $K_0 \in \mathcal{K}$, $N_{\mathrm{traj}}$, $l$, $r$, fixed step size $\eta$, weight $\lambda$,
        tolerance $\varepsilon$
        \FOR{$j = 0, 1, 2, \dots$}
            \STATE compute $\widehat{\nabla f}(K_j)$ by \cref{alg:mf-gradient}
            \STATE compute $K_{j+1}$ by \eqref{eq:spi-prox}, with $\widehat{\nabla f}(K_j)$ in place of $\nabla f(K_j)$
            \IF{$\lVert K_{j+1} - K_j \rVert_F \leq \varepsilon$}
                \STATE \textbf{stop}
            \ENDIF
        \ENDFOR
    \end{algorithmic}
\end{algorithm}

Now to the selection of the parameters. Park et al.\ already note the difficulty of finding a fixed step size
that does not produce divergence of the method~\cite{Park2020Structured}. Also, the example shown only
uses the model-based approach and the regularizer to induce sparse structure, not the model-free variant combined
with low dimensionality. So during our first tests, we often ran into either unstable controllers, or ones that
did not induce any structure at all, hinting at a too low $\lambda$. If we remember the proximal step
\eqref{eq:spi-prox}, we can see that it depends on both the step size and $\lambda$, as every singular value is
mapped to
\begin{equation}
    \label{eq:spi-threshold}
    \sigma_i \mapsto \max\{ \sigma_i - \eta \lambda, 0 \} .
\end{equation}
Even after some improvement by using the two-point estimator (\cref{alg:mf-gradient}), it remained quite a difficult
topic. So rather extensive experiments on the variation of $\lambda$ and the step size can be found in \cref{sec:app-spi-mf}.
Based on those, $\lambda = 4000$ and a step size of $\eta = 3.5 \times 10^{-6}$ were selected.

\subsection{Results}
\label{subsec:unknown-model-results}

\subsection{Noise Experiments on Large Inflation}
\label{subsec:noise-large-inflation}

\subsection{Discussion}
\label{sec:id-model-discussion}

\section{Discussion}
\label{sec:ch4-discussion}
